\documentclass{article}
\usepackage{graphicx}
\usepackage[margin=1in]{geometry}
\usepackage{enumitem}
\usepackage{booktabs}
\usepackage{array}
\usepackage{parskip}
\usepackage{hyperref}
\setlist{noitemsep, topsep=3pt, leftmargin=*}

\title{Beyond the Shadow of a Doubt: Benchmarking MCMC Sampling Algorithms
  for Multimodal Astrophysical Inference\\[0.5em]
  \large Progress Report}
\author{Lucas Arthur \and Samson Mercier}
\date{April 13, 2026}

\begin{document}

\maketitle

\begin{abstract}
As astrophysics enters the era of petabyte-scale datasets, the choice of
inference algorithm and implementation for a research problem could make the
difference between an intractable or incorrect analysis and a major scientific
discovery. The high dimensionality and multi-modality of posterior distributions
can pose problems for popular inference tools. To ensure that the astrophysics
community is fully leveraging available inference algorithms, we benchmark a
set of widespread tools from across the astrophysical literature. We benchmark
eight inference algorithms---random walk Metropolis-Hastings (RWMH),
affine-invariant ensemble sampling, parallel tempering, NUTS, SMC with
tempering, non-reversible parallel tempering, ensemble Kalman sampling, and
nested sampling (as a non-MCMC comparator)---on two test problems: a synthetic
8-component 2D Gaussian mixture with known ground truth, and a simulated
exoplanet transit with stellar spots. All experiments run on AMD EPYC 9474F
48-core CPUs and Nvidia L40S GPUs using JAX to ensure
implementation-agnostic comparison. We evaluate each sampler on mode
recovery, posterior accuracy, and computational efficiency (ESS per unit compute
time), and will produce a concrete set of recommendations for practitioners.
\end{abstract}

% -----------------------------------------------------------------------
\section{Project Overview}

The central question of this project is: \textit{Which parallelizable
MCMC-family samplers most reliably recover multimodal posteriors in
astrophysical settings, and at what computational cost?} We benchmark eight
algorithms: RWMH, affine-invariant ensemble sampling \cite{goodman2010,emcee},
reversible parallel tempering (SEO) \cite{vousden2016}, non-reversible parallel
tempering (DEO) \cite{syed2019}, NUTS, SMC with tempering \cite{delmoral2006},
ensemble Kalman sampling \cite{garbuno2020}, and nested sampling
\cite{speagle2020} as a non-MCMC reference.

We evaluate each sampler on two test problems. The first is an 8-component 2D
Gaussian mixture with fully known ground-truth modes, weights, and covariances,
enabling precise measurement of mode recovery and posterior accuracy. The second
is a simulated exoplanet transit with stellar spots, a problem known to produce
multimodal posteriors, which tests each sampler in a realistic astrophysical
setting.

The primary evaluation criteria are: mode recovery (fraction of modes recovered,
recovered mode weights), posterior accuracy (Wasserstein-2 distance from ground
truth where available), and computational efficiency (ESS per log-density
evaluation and ESS per wall-clock second). Secondary diagnostics
include $\hat{R}$, acceptance rates, and---for parallel tempering variants---swap
acceptance rates and round-trip counts between temperature levels.

The full implementation stack is JAX~\cite{jax2018github} / BlackJAX~\cite{blackjax} /
NumPyro~\cite{numpyro}, ensuring hardware-accelerated, implementation-agnostic
comparison across all algorithms.

% -----------------------------------------------------------------------
\section{Division of Labor Update}

The division of labor has been reshuffled relative to the original proposal.
Lucas Arthur has taken ownership of both parallel tempering variants---reversible
SEO and non-reversible DEO---in addition to his originally assigned algorithms
(NUTS and SMC with tempering). This consolidation arose naturally because the two
tempering variants share substantial infrastructure (temperature ladder
management, swap kernels, diagnostic logging), and keeping them under a single
owner reduced integration overhead.

Samson Mercier has taken on all forward model development, including
implementation and validation of the stellar spot transit model against external
transit-fitting codes, in addition to his originally assigned algorithms (RWMH,
affine-invariant MCMC, and nested sampling).

The ensemble Kalman sampler (EKS) is deprioritized. Implementation is on hold
until all other samplers have been run and results collated across both test
problems. Samson Mercier will implement it if time permits; it will be omitted
from the final comparison if it is not ready in time.

% -----------------------------------------------------------------------
\section{Progress to Date}

\subsection{Gaussian Mixture Model --- Substantially Complete}

Seven of the eight samplers have been fully implemented and run on the Gaussian
mixture problem: RWMH, affine-invariant MCMC, reversible PT (SEO),
non-reversible PT (DEO), NUTS, SMC with tempering, and nested sampling. EKS is
the sole missing algorithm; it is deferred pending completion of transit model
runs.

Each sampler produces the following outputs: posterior samples in a
standardized format, corner plots, and a diagnostics record containing ESS,
$\hat{R}$, acceptance rates, recovered mode weights, and inter-mode transition
counts (for chain-based methods).

\subsection{Stellar Spot Transit Model --- Forward Model Complete, Sampling Not
Yet Started}

The transit forward model has been implemented using the \texttt{sajax} library
with a NumPyro likelihood and a BlackJAX-compatible log-density interface. The
inference problem is to recover three spot parameters---latitude, longitude, and
radius---from a synthetic noisy light curve generated from a known ground-truth
configuration. This is a reduced scope relative to the proposal, which described
a 12-parameter model; the three-parameter version was chosen to keep the transit
problem tractable while preserving the multimodal structure of interest. Samson
Mercier has validated the forward model against independent transit-fitting codes.

Sampler scripts for the transit problem have not yet been written. Writing and
running these scripts is the primary remaining task.

% -----------------------------------------------------------------------
\section{Preliminary Results: Gaussian Mixture}

Detailed quantitative analysis is reserved for the final report; we summarize
high-level qualitative observations here. Notable observations include:

\begin{itemize}
  \item \textbf{RWMH} shows poor mixing between modes. The chain rarely crosses
    the low-density region separating mixture components, resulting in
    substantially fewer inter-mode transitions and lower ESS per evaluation than
    all other methods.
  \item \textbf{Non-reversible PT (DEO)} achieves substantially more inter-mode
    transitions than reversible PT (SEO) at matched compute budgets, consistent
    with the theoretical prediction that non-reversible temperature swaps improve
    round-trip rates \cite{syed2019}.
  \item \textbf{Nested sampling} reliably recovers all eight modes. Its ESS per
    log-density evaluation is lower than gradient-based methods, but it provides
    an accurate evidence estimate and serves as a robust reference.
  \item \textbf{NUTS} and \textbf{SMC with tempering} show strong ESS per
    evaluation on this low-dimensional problem; their relative performance on the
    higher-dimensional transit model remains to be seen.
\end{itemize}

% -----------------------------------------------------------------------
\section{Revised Plan and Remaining Tasks}

Table~\ref{tab:plan} lists the remaining tasks and internal deadlines.

\begin{table}[h]
\centering\small
\begin{tabular}{@{}p{0.60\textwidth}p{0.14\textwidth}p{0.10\textwidth}@{}}
\toprule
\textbf{Task} & \textbf{Owner} & \textbf{Target} \\
\midrule
Run all 7 implemented samplers on transit model & Both   & Apr 25 \\
Collate all results; build comparison tables    & Both   & Apr 25 \\
Implement EKS (if time permits)                 & Samson & Apr 25 \\
Produce final figures                           & Both   & Apr 28 \\
Write final report                              & Both   & May 1  \\
\bottomrule
\end{tabular}
\caption{Remaining tasks and internal deadlines.}
\label{tab:plan}
\end{table}

% -----------------------------------------------------------------------
\section{Planned Figures and Evaluation Framework}

Comparing the eight algorithms requires care because they produce
fundamentally different outputs: correlated samples (RWMH, PT, NUTS), weighted
particle sets (SMC, nested sampling), or continuously evolved ensembles (EKS,
affine-invariant). We therefore report efficiency as ESS per log-likelihood
evaluation, the most universally applicable unit of compute. For chain-based
methods, warm-up evaluations are included in the denominator to give an honest
accounting of total compute. The primary mode-recovery metrics are the per-mode
weight error $|\hat{w}_k - w_k|$, total variation distance from the true
mixture, and (for chain methods) inter-mode transition count. For the Gaussian
mixture we also report Wasserstein-2 distance to ground truth; for the transit
model we use nested-sampling output as a reference distribution.

Anticipated figures: (1) KDE posterior overlays for all samplers on the 2D
mixture against the true density; (2) per-mode weight recovery bar charts;
(3) ESS per likelihood evaluation across samplers and both problems; (4)
convergence curves (Wasserstein distance vs.\ evaluations); (5) PT swap
acceptance and round-trip rates for SEO vs.\ DEO \cite{syed2019}; (6) transit
model corner plots per sampler; (7) a summary table of $\hat{R}$, MCSE,
ESS/eval, and modes recovered across both problems.

\bibliographystyle{plain}
\bibliography{refs}

\end{document}
